\documentclass[11pt,a4paper]{article}
\usepackage[UTF8,fontset=fandol]{ctex}
\usepackage{amsmath,amssymb}
\usepackage{booktabs}
\usepackage{multirow}
\usepackage{graphicx}
\usepackage{geometry}
\geometry{left=2.0cm,right=2.0cm,top=2.2cm,bottom=2.2cm}
\usepackage{caption}
\captionsetup{font=small,labelfont=bf}
\usepackage{longtable}
\usepackage{array}

\newcommand{\suppdir}{../tables_v3}

\title{\textbf{补充材料：支持表格 S1--S9}\\[0.4em]
\large Supplementary Tables for\\
\emph{Mechanistic Benchmarks and Empirical Calibration of Epidemic Forecast Horizons}}
\author{（双盲评审版，作者信息暂略）}
\date{\today}

\begin{document}
\maketitle

\noindent\textbf{说明。}本补充材料收录正文引用的全部支持表格（表 S1--S9）。所有表格均由随稿提交的代码与数据冻结快照自动生成：表 S2 来自 \texttt{reports/verify\_t4.json}，表 S5 来自 \texttt{reports\_v3/caliber\_table.json}，表 S8 来自 \texttt{reports\_v3/flusight\_v1.*\_extended.json}，其余表格由 \texttt{scripts\_v2/} 下的生成脚本产出并受 \texttt{scripts\_v2/check\_consistency.py} 的一致性门禁逐字节核验。

\section*{表 S1　主要符号与口径对照表}
\addcontentsline{toc}{section}{表 S1 主要符号}
\begin{table}[htbp]\centering\footnotesize
\caption{全文主要符号、名称与定义口径}
\begin{tabular}{@{}>{\raggedright\arraybackslash}p{2.6cm}>{\raggedright\arraybackslash}p{3.2cm}>{\raggedright\arraybackslash}p{9.0cm}@{}}
\toprule
符号 & 名称 & 定义与口径 \\
\midrule
$\mathcal{F}_t$ & 预测原点信息集 & 原点 $t$ 处严格可得的全部历史与外生信息 \\
$Y_{t+h}$ & 预测目标 & 前瞻 $h$ 步的真实状态（终报真值口径下 $Y_{t+h}=Z_{t+h}$） \\
$\delta_h$ & 预测规则 & $\mathcal{F}_t$ 可测且二阶可积的实随机变量 \\
$m_h,\,V_h$ & 条件均值与条件方差 & $m_h=\mathbb{E}[Y_{t+h}\mid\mathcal{F}_t]$，$V_h=\operatorname{Var}(Y_{t+h}\mid\mathcal{F}_t)$ \\
$\mathcal{R}_h(\delta_h)$ & 条件平方预测风险 & $\mathbb{E}[(Y_{t+h}-\delta_h)^2\mid\mathcal{F}_t]$；最小值为地板 $V_h$ \\
$Z_t$ & 周度宏观观测 & 第 $t$ 周全辖区报告的新增住院计数 \\
$R$ & 代际增长乘子 & 个体级（代际时钟）基本再生数 \\
$R_{\text{周}}$ & 周度增长乘子 & $\exp(b_{\text{week}})$；宏观主口径使用 \\
$s_{\text{周}}$ & 斜率对数标准误 & $\operatorname{SE}(b_{\text{week}})$，自由度 $3$（$W=5$） \\
$k_{\text{ind}}$ & 个体级离散度 & 个体子代分布负二项的 size 参数 \\
$k_{\text{agg}}$ & 宏观聚合离散度 & 州级周度序列的有效离散度（对数差分二阶矩估计，上限 1,000） \\
$I_0$ & 初始发病规模 & 推断窗口内周度均值（平滑初始状态） \\
$\text{CV}^2_{\text{macro}}(h)$ & 宏观过程相对方差 & 固定 $k_{\text{agg}}$ 的 NB2 更新模型有限步精确方差除以均值平方 \\
$P_{\text{exact}}(h)$ & 参数外推项（精确） & $\exp(2h^2s^2)-2\exp(h^2s^2/2)+1$（对数正态精确式） \\
$P_{\text{quad}}(h)$ & 参数外推项（领先阶） & $(h_{\text{周}}s_{\text{周}})^2$，用于误差记账 \\
$h^*$ & 机制视界 & $\text{RelMSE}(h)$ 首次达到容忍度 $\tau$ 的连续阈值根 \\
$h^*_{\text{obs}}$ & 实测绝对误差视界 & 各州滚动误差首次穿越同一 $	au$ 的插值点 \\
$\tau$ & 容忍度阈值 & 本文主口径 $\tau=0.5$ \\
$\mathcal{E}_{\text{res}}(h)$ & 模型—观测代数差额 & $\text{RelMSE}_{\text{obs}}-(\text{CV}^2_{\text{macro}}+P_{\text{quad}})$，描述性记账 \\
$\phi,\,s_e^2$ & AR(1) 持续性参数 & 环境漂移自相关系数与其创新方差（$|\phi|<1$） \\
$\text{WIS}$ & 加权区间评分 & 基于 23 个标准分位数的严格适当评分规则 \\
$\text{PIT}$ & 概率积分变换 & 分位数插值近似随机化 PIT，完美校准时服从均匀分布 \\
$C$ & 理想样本量基准 & 给定效应量与功效下由 CRB 反解的单代完全观测样本量 \\
\bottomrule
\end{tabular}

\end{table}

\clearpage
\section*{表 S2　理想样本量基准与蒙特卡洛参数全网格}
\addcontentsline{toc}{section}{表 S2 样本量基准与模拟网格}
\begin{table}[htbp]\centering\footnotesize
\caption{定理 4（Fisher 信息与 Cramér--Rao 界）验证的完整参数网格。$C$ 为单代完全观测的样本量基准（单位信噪比 $C=861$、80\% 功效工作点 $C=1{,}300$，以及本文实测 $k$ 对应的 $C=2{,}162$ 与 $C=6{,}593$）；经验方差为 50{,}000 次蒙特卡洛数值极大似然的方差}
\begin{tabular}{ccccrrrr}
\toprule
$R$ & $k$ & $C$ & 重复次数 & 经验方差 & CRB & 经验/CRB & 偏差 \\
\midrule
1.05 & 1.000 & 200 & 50000 & 1.084e-02 & 1.076e-02 & 1.0073 & +6.57e-04 \\
1.05 & 1.000 & 1300 & 50000 & 1.674e-03 & 1.656e-03 & 1.0112 & +1.59e-05 \\
1.10 & 1.000 & 861 & 50000 & 2.700e-03 & 2.683e-03 & 1.0062 & -8.71e-05 \\
1.10 & 1.000 & 1300 & 50000 & 1.770e-03 & 1.777e-03 & 0.9964 & +9.47e-05 \\
1.10 & 0.434 & 2162 & 50000 & 1.789e-03 & 1.798e-03 & 0.9950 & +5.64e-05 \\
1.10 & 0.111 & 6593 & 50000 & 1.805e-03 & 1.820e-03 & 0.9917 & +1.49e-04 \\
1.30 & 1.000 & 500 & 50000 & 5.968e-03 & 5.980e-03 & 0.9980 & +2.02e-05 \\
1.30 & 1.000 & 2000 & 50000 & 1.510e-03 & 1.495e-03 & 1.0100 & +2.06e-04 \\
1.50 & 0.300 & 500 & 50000 & 1.811e-02 & 1.800e-02 & 1.0060 & +7.88e-04 \\
1.50 & 0.300 & 2000 & 50000 & 4.465e-03 & 4.500e-03 & 0.9923 & -1.48e-04 \\
\bottomrule
\end{tabular}

\end{table}

\section*{表 S3　各分析对象的数据来源与潜在泄漏核查表}
\addcontentsline{toc}{section}{表 S3 数据源与泄漏核查}
\begin{table}[htbp]\centering\footnotesize
\caption{各分析对象的参数来源、真值口径与未来信息依赖}
\begin{tabular}{@{}>{\raggedright\arraybackslash}p{2.5cm}>{\raggedright\arraybackslash}p{2.4cm}>{\raggedright\arraybackslash}p{2.2cm}>{\raggedright\arraybackslash}p{2.6cm}>{\raggedright\arraybackslash}p{3.0cm}@{}}
\toprule
分析对象 & 参数来源 & 真值来源 & 是否使用未来信息（泄漏风险） & 结论性质 \\
\midrule
理论模拟验证 & 已知真参数 & 模拟轨迹 & 否 & 严格数值验证 \\
微观传播链队列 & 完整传播链（回顾） & 回顾性传播链 & 可能存在追踪截断与右删失 & 机制拟合与模型内自洽性 \\
州级条件机制视界 & 5 周滑动窗口（在线可得） & 冻结面板（Final Vintage） & 参数可能受后续回填修订影响 & 回溯性机制估计 \\
动力学阶段分类 & 全季轨迹曲率与峰值 & 全季最终真值 & \textbf{是}，依赖未来峰值 & 仅为事后分层，非实时可用 \\
FluSight 外部审计 & 官方提交分位数（原点可得） & Final Vintage Truth & 真值含事后回填 & 外部回溯审计 \\
逐州滚动回测 & 历史原点前 4 周参数 & Final Vintage Truth & 真值版本含未来回填 & 终报口径业务时效 \\
\bottomrule
\end{tabular}
\end{table}

\section*{表 S4　四类视界口径的定义与层级归属}
\addcontentsline{toc}{section}{表 S4 视界口径定义}
\begin{table}[htbp]\centering\footnotesize
\caption{四类视界的生成模型、层级归属与适用数据}
\begin{tabular}{@{}>{\raggedright\arraybackslash}p{3.4cm}>{\raggedright\arraybackslash}p{4.0cm}>{\raggedright\arraybackslash}p{2.6cm}>{\raggedright\arraybackslash}p{3.2cm}@{}}
\toprule
视界名称 & 生成模型 & 层级归属 & 适用数据 \\
\midrule
微观机制视界 $h^*_{\text{micro}}$ & Galton--Watson 负二项分支（子代方差线性于 $Z_{t-1}$） & 理论主结果（个体级） & 个体传播链队列 \\
启发式尺度映射视界 $h^*_{\text{heur}}$ & 微观分支闭式代入宏观离散度 $k_{\text{agg}}$（代际时钟，模型不自洽） & 敏感性分析 & 州级序列（仅作对照） \\
宏观条件精确机制视界 $h^*_{\text{macro}}$ & 固定 $k_{\text{agg}}$ 负二项宏观更新模型（条件方差含 $Z_{t-1}^2$，周度时钟） & 州级主结果 & 周度住院面板 \\
实测业务视界 $h^*_{\text{obs}}$ & 观测滚动误差首次穿越容忍度 $\tau$ & 现实比较 & 州级预测回测 \\
\bottomrule
\end{tabular}
\end{table}

\clearpage
\section*{表 S5　七阶段六口径机制视界全景对照}
\addcontentsline{toc}{section}{表 S5 六口径视界全景}
\begin{table}[htbp]\centering\scriptsize
\caption{各阶段在六类机制视界口径下的州级中位视界（周），并与同定义实测绝对误差视界及相对持续性基线的业务时效对照。$h^*_{\text{macro}}$ 为正文主口径（固定 $k_{\text{agg}}$ 的 NB2 更新模型、周度时钟）；$h^*_{k}$ 三列为代入三档个体级离散度（实测 $\hat k$）的启发式尺度映射，仅作平行敏感性对照}
\begin{tabular}{@{}lrrrrrrrrr@{}}
\toprule
阶段 & $n$ & $h^*_{\text{macro}}$ & $h^*_{\text{macro,lin}}$ & $h^*_{\text{heur}}$ & $h^*_{k=0.434}$ & $h^*_{k=0.182}$ & $h^*_{k=0.111}$ & $h^*_{\text{obs}}$ & $h^*_{\text{persist}}$ \\
\midrule
COVID-19 Delta 暴发期 & 50 & 7.12 & 8.17 & 10.86 & 9.07 & 8.12 & 5.91 & 1.78 & 3.72 \\
COVID-19 Omicron 达峰期 & 51 & 3.95 & 4.15 & 8.02 & 6.69 & 5.83 & 4.89 & 1.92 & 1.43 \\
COVID-19 JN.1 流行期 & 51 & 5.00 & 5.46 & 8.69 & 7.11 & 5.59 & 4.48 & 3.89 & 1.00 \\
流感 2022--23 暴发早期 & 38 & 2.20 & 2.46 & 4.47 & 2.29 & 1.10 & 0.67 & 1.23 & 2.80 \\
流感 2024--25 流行季 & 45 & 2.58 & 3.13 & 6.15 & 4.81 & 2.45 & 1.30 & 1.27 & 2.86 \\
RSV 2024--25 流行季 & 34 & 5.82 & 6.41 & 9.60 & 6.75 & 4.50 & 2.85 & 2.91 & 1.28 \\
RSV 2025--26 流行季 & 26 & 3.89 & 4.79 & 6.74 & 4.11 & 1.92 & 1.15 & 3.49 & 4.26 \\
\bottomrule
\end{tabular}

\end{table}

\section*{表 S6　推断窗口与 $k_{\text{agg}}$ 截断敏感性矩阵}
\addcontentsline{toc}{section}{表 S6 窗口与截断敏感性}
\begin{table}[htbp]\centering\footnotesize
\caption{各阶段州级中位机制视界对推断窗口长度（5/6/8/10 周）与 $k_{\text{agg}}$ 上限截断（100/500/1{,}000/无约束）的敏感性}
\begin{tabular}{lrrrrrrrrr}
\toprule
& \multicolumn{4}{c}{推断窗口长度（周）} & \multicolumn{4}{c}{$k_{\text{agg}}$ 上限截断} \\
\cmidrule(lr){2-5}\cmidrule(lr){6-9}
阶段 & 5 & 6 & 8 & 10 & 100 & 500 & 1000 & $\infty$ \\
\midrule
COVID Delta        & 7.12 & 7.50 & 8.07 & 6.79 & 7.12 & 7.12 & 7.12 & 7.12 \\
COVID Omicron      & 3.95 & 4.15 & 4.17 & 3.61 & 3.95 & 3.95 & 3.95 & 3.95 \\
COVID JN.1         & 5.00 & 4.24 & 5.77 & 7.23 & 5.00 & 5.00 & 5.00 & 5.00 \\
Influenza 2022--23 & 2.20 & 2.39 & 2.00 & 1.79 & 2.13 & 2.19 & 2.20 & 2.20 \\
Influenza 2024--25 & 2.58 & 3.03 & 1.52 & 1.67 & 2.58 & 2.58 & 2.58 & 2.58 \\
RSV 2024--25       & 5.82 & 5.57 & 6.21 & 2.98 & 5.54 & 5.82 & 5.82 & 5.82 \\
RSV 2025--26       & 3.89 & 4.82 & 5.58 & 4.79 & 3.60 & 3.85 & 3.89 & 3.89 \\
\bottomrule
\end{tabular}

\end{table}

\section*{表 S7　不同容错门槛下的多层级预警映射}
\addcontentsline{toc}{section}{表 S7 容错门槛预警映射}
\begin{table}[htbp]\centering\footnotesize
\caption{各阶段州级机制视界在四种容忍度门槛 $\tau$ 下的取值（括号内为参与辖区数）}
\begin{tabular}{lcccc}
\toprule
阶段 & $\tau=0.20$ & $\tau=0.35$ & $\tau=0.50$ & $\tau=0.70$ \\
\midrule
COVID-19 Delta 暴发期 & 1.5 (50) & 4.1 (50) & 7.1 (50) & 11.0 (50) \\
COVID-19 Omicron 达峰期 & 0.8 (51) & 2.2 (51) & 4.0 (51) & 6.4 (51) \\
COVID-19 JN.1 流行期 & 1.0 (51) & 2.8 (51) & 5.0 (51) & 7.9 (51) \\
流感 2022--23 暴发早期 & 0.4 (38) & 1.2 (38) & 2.2 (38) & 3.9 (38) \\
流感 2024--25 流行季 & 0.5 (45) & 1.4 (45) & 2.6 (45) & 4.4 (45) \\
RSV 2024--25 流行季 & 1.0 (34) & 3.3 (34) & 5.8 (34) & 9.1 (34) \\
RSV 2025--26 流行季 & 0.6 (26) & 1.9 (26) & 3.9 (26) & 6.7 (26) \\
\bottomrule
\end{tabular}

\end{table}

\clearpage
\section*{表 S8　CDC FluSight 三赛季分步长完整评分面板}
\addcontentsline{toc}{section}{表 S8 FluSight 分步长评分面板}
\begin{table}[htbp]\centering\scriptsize
\caption{三个锁定赛季在 $h=1,2,3$ 周提前期下的三模型完整评分面板。每行三个模型在同一（赛季, 提前期）单元内共享完全相同的 $n$（严格四方交集），故差值可直接读作配对性能差。逐单元、逐辖区的明细见随包 \texttt{reports\_v3/flusight\_v1.*.csv}}
\begin{tabular}{@{}llrrrrrrrr@{}}
\toprule
赛季 & $h$（周） & $n$ & WIS & 覆盖率 50\% & 覆盖率 80\% & 覆盖率 95\% & 80\% 区间宽 & 95\% 区间宽 & PIT 均值 \\
\midrule
\multirow{3}{*}{v1.0.0 (2023--24)} & \multirow{3}{*}{1} & 1,337 & 34.00 & 0.524 & 0.794 & 0.933 & 150.8 & 231.8 & 0.568 \\
& & 1,337 & 45.68 & 0.265 & 0.606 & 0.895 & 106.1 & 256.5 & 0.500 \\
& & 1,337 & 59.94 & 0.214 & 0.384 & 0.554 & 92.3 & 142.1 & 0.506 \\
\midrule
\multirow{3}{*}{v1.0.0 (2023--24)} & \multirow{3}{*}{2} & 1,286 & 45.55 & 0.474 & 0.765 & 0.925 & 185.5 & 286.9 & 0.580 \\
& & 1,286 & 62.63 & 0.250 & 0.608 & 0.863 & 146.7 & 290.2 & 0.500 \\
& & 1,286 & 68.63 & 0.252 & 0.472 & 0.664 & 152.1 & 237.5 & 0.508 \\
\midrule
\multirow{3}{*}{v1.0.0 (2023--24)} & \multirow{3}{*}{3} & 1,235 & 54.60 & 0.478 & 0.737 & 0.901 & 213.7 & 332.4 & 0.590 \\
& & 1,235 & 76.72 & 0.247 & 0.608 & 0.854 & 178.0 & 318.1 & 0.499 \\
& & 1,235 & 87.98 & 0.271 & 0.517 & 0.714 & 227.3 & 363.4 & 0.509 \\
\midrule
\multirow{3}{*}{v1.1.0 (2024--25)} & \multirow{3}{*}{1} & 1,315 & 85.27 & 0.551 & 0.739 & 0.837 & 255.9 & 382.1 & 0.605 \\
& & 1,315 & 125.11 & 0.317 & 0.552 & 0.742 & 193.1 & 376.4 & 0.479 \\
& & 1,315 & 170.88 & 0.128 & 0.269 & 0.405 & 213.2 & 329.0 & 0.381 \\
\midrule
\multirow{3}{*}{v1.1.0 (2024--25)} & \multirow{3}{*}{2} & 1,315 & 110.33 & 0.525 & 0.710 & 0.820 & 282.4 & 426.6 & 0.603 \\
& & 1,315 & 170.89 & 0.330 & 0.525 & 0.690 & 248.2 & 430.8 & 0.466 \\
& & 1,315 & 181.61 & 0.221 & 0.404 & 0.548 & 362.4 & 571.5 & 0.366 \\
\midrule
\multirow{3}{*}{v1.1.0 (2024--25)} & \multirow{3}{*}{3} & 1,315 & 138.62 & 0.501 & 0.695 & 0.801 & 290.9 & 441.7 & 0.607 \\
& & 1,315 & 216.96 & 0.342 & 0.500 & 0.650 & 288.6 & 476.8 & 0.462 \\
& & 1,315 & 222.56 & 0.284 & 0.471 & 0.615 & 568.3 & 931.3 & 0.365 \\
\midrule
\multirow{3}{*}{v1.2.0 (2025--26)} & \multirow{3}{*}{1} & 1,376 & 60.24 & 0.516 & 0.780 & 0.905 & 199.9 & 315.1 & 0.559 \\
& & 1,376 & 90.57 & 0.438 & 0.738 & 0.869 & 204.7 & 407.1 & 0.486 \\
& & 1,376 & 116.72 & 0.201 & 0.378 & 0.531 & 174.2 & 269.7 & 0.363 \\
\midrule
\multirow{3}{*}{v1.2.0 (2025--26)} & \multirow{3}{*}{2} & 1,376 & 79.97 & 0.509 & 0.764 & 0.890 & 237.8 & 377.3 & 0.541 \\
& & 1,376 & 116.77 & 0.477 & 0.739 & 0.853 & 266.4 & 463.2 & 0.464 \\
& & 1,376 & 141.23 & 0.311 & 0.493 & 0.649 & 300.9 & 477.7 & 0.389 \\
\midrule
\multirow{3}{*}{v1.2.0 (2025--26)} & \multirow{3}{*}{3} & 1,376 & 91.15 & 0.520 & 0.762 & 0.880 & 263.3 & 415.8 & 0.520 \\
& & 1,376 & 133.49 & 0.493 & 0.730 & 0.839 & 310.2 & 507.6 & 0.444 \\
& & 1,376 & 191.55 & 0.377 & 0.564 & 0.690 & 489.5 & 809.2 & 0.412 \\
\bottomrule
\end{tabular}

\end{table}

\section*{表 S9　COVIDhub-ensemble 12 预测原点逐周技能衰减明细}
\addcontentsline{toc}{section}{表 S9 COVIDhub 逐原点明细}
\begin{table}[htbp]\centering\scriptsize
\caption{Delta 与 Omicron 两波共 12 个预测原点、$h=1$--$4$ 周的跨州中位平方误差比率（相对持续性基线）与绝对加权区间评分明细}
\input{\suppdir/table6_hub}
\end{table}

\end{document}
